\section{Conclusion}


We have presented \ConservAI{}, an integrated deep learning pipeline for wildlife population monitoring and habitat quality assessment. By combining hierarchical species classification, temporal context attention, automated population estimation, and multi-spectral habitat analysis, \ConservAI{} provides conservation practitioners with a comprehensive, automated monitoring system. Deployment on the \Zoo{} decentralized infrastructure ensures data sovereignty, model provenance, and reproducible evaluation. Our results on four benchmark datasets and a 12-month field deployment demonstrate that AI-powered monitoring can match or exceed the accuracy of traditional methods while reducing processing time by over 200$\times$.

The \Zoo{} Labs Foundation releases all models, training code, and evaluation protocols as open-source contributions to the conservation AI community. We invite conservation organizations to join the \Zoo{} network and contribute camera trap data to the Data Commons, enabling the development of increasingly accurate and comprehensive monitoring tools.

